\section{Analysis}

\subsection{Contribution of Each Component}

\textbf{Table 8.1: Ablation of components (DistilBERT).}

\begin{table}[H]
\centering
\small
\begin{tabular}{lll}
\toprule
\textbf{Config} & \textbf{Accuracy} & \textbf{$\Delta$} \\
\midrule
Baseline (linear attention only) & {{D\_BASE\_ATTN}} & -- \\
+ Dual-scale RFF & {{D\_BASE\_DUAL}} & +{{DUAL\_GAIN}} \\
+ Landmark distillation & {{D\_BASE\_LAND}} & +{{LANDMARK\_GAIN}} \\
+ Local window & {{D\_BASE\_LOCAL}} & +{{LOCAL\_GAIN}} \\
+ Softpick & \textbf{{{D\_ACC}}} & +{{SOFTPICK\_GAIN}} \\
\bottomrule
\end{tabular}
\end{table}

All four components contribute.

\subsection{Landmark Selection Analysis}

\textbf{Figure 8.1} shows the distribution of selected landmarks.

\begin{itemize}
\item \textbf{Position 0} is always retained (guaranteed).
\item \textbf{Token 1-5} are frequently selected (early context).
\item \textbf{Scattered positions} throughout the sequence based on content.

\end{itemize}
\subsection{Why Softpick Matters}

\textbf{Figure 8.2} shows the case where a full row of attention is masked. 
Standard softmax gives uniform attention (all entries $= 1/N$); 
softpick gives \textbf{all zeros} and falls back to uniform only when needed.

\textbf{Result.} Softpick prevents numerical instability at masked rows.

\subsection{Local Window Contribution}

\textbf{Table 8.2: Effect of window size.}

\begin{table}[H]
\centering
\small
\begin{tabular}{lll}
\toprule
\textbf{W} & \textbf{DistilBERT} & \textbf{Efficiency} \\
\midrule
32 & {{D\_W32}} & Fastest \\
{{LOCAL\_WINDOW}} & {{D\_ACC}} & Sweet spot \\
256 & {{D\_W256}} & Slower \\
Full L & {{D\_WFULL}} & Same as MHSA \\
\bottomrule
\end{tabular}
\end{table}

\subsection{Dual-Scale Analysis}

The two RFF projections capture different frequency bands. \textbf{Figure 8.3} 
shows the singular value spectra of $W_1$ and $W_2$:

\begin{itemize}
\item $W_1$ (F=128): captures low-frequency kernel features.
\item $W_2$ (F/2=64): captures high-frequency kernel features.

\end{itemize}
Combining them gives a richer representation.

\subsection{Learned Local Gate}

\textbf{Figure 8.4} shows per-head gate values after training. Average values:

\begin{table}[H]
\centering
\small
\begin{tabular}{ll}
\toprule
\textbf{Model} & \textbf{Mean $g_h$} \\
\midrule
DistilBERT & {{D\_G\_MEAN}} \\
RoBERTa & {{R\_G\_MEAN}} \\
GPT-2 & {{G\_G\_MEAN}} \\
\bottomrule
\end{tabular}
\end{table}

Heads with high $g_h$ prefer global; low $g_h$ prefer local.

\subsection{Complexity Reality Check}

\textbf{Table 8.3: Actual complexity.}

\begin{table}[H]
\centering
\small
\begin{tabular}{ll}
\toprule
\textbf{Component} & \textbf{Complexity} \\
\midrule
Local window & $O(L \cdot W)$ \\
Global landmark & $O(L \cdot L_\text{max})$ \\
\textbf{Total} & \textbf{$O(L \cdot (W + L_\text{max}))$} \\
\bottomrule
\end{tabular}
\end{table}

With $L_\text{max} = \rho L$ and $\rho = 0.5$, the landmark term is 
$O(L^2 / 2)$ -- still quadratic but with half the constant.

\textbf{Note.} LEKAv2 is \textbf{not truly linear} unless $L_\text{max}$ is fixed.

\subsection{Failure Modes}

\subsubsection{Fatal Error Handling}

The \texttt{copy\_weights} still raises \texttt{RuntimeError} on extractor failure.

\subsubsection{Landmark Saturation}

If the scorer assigns equal scores to all tokens, selection is effectively 
random.

\subsubsection{Local Gate Saturation}

If $g_h \to 0$ or $g_h \to 1$, one branch is ignored.

\subsection{Summary of Insights}

\begin{enumerate}
\item \textbf{All four components contribute.}
\item \textbf{Landmark distillation is essential.}
\item \textbf{Dual-scale helps.}
\item \textbf{Softpick stabilizes.}
\item \textbf{LEKAv2 is not truly linear at   = 0.5.}
\end{enumerate}
\begin{lstlisting}

---
\end{lstlisting}